\section{Discovery and Search}
\label{sec:discovery}

\subsection{Search on the device}

Search runs over listing headers already held by the client: a case-insensitive substring match over title, category, ship-from region, kind, description, and seller name. The query is not sent anywhere, so neither the node, a replica, nor an index learns the search terms. The cost is that the client must hold the headers it searches. We measured this filter at 0.2~ms per query over 1{,}000 headers and 5~ms over 10{,}000 headers on a laptop (\S\ref{sec:evaluation}); 100{,}000 headers took 49--67~ms per query. Matching is literal; there is no ranking, stemming, or index. Results then pass through the buyer's content filter (\S\ref{sec:labels}), which also runs on the device.

\subsection{Finding stores and nodes}

The design treats every source of nodes as untrusted for content and only partly trusted for availability:
\begin{itemize}
  \item A user can paste a store link (a store key plus an onion address). A store reached this way opens with an empty seed list.
  \item Seeds shipped with the software, and peers remembered from earlier sessions, are used to find replicas.
  \item Every record received from any of them is verified against its signature and content hash before use.
\end{itemize}

Because records are self-authenticating, a malicious seed or replica cannot forge or alter a listing. It can withhold records, serve stale ones (\S\ref{sec:catalog}), or, if it controls all of a node's peers, present a partial view of the network: an eclipse attack~\citep{heilman2015eclipse}. Remembering peers across sessions, drawing peers from several independent sources, and preferring a pasted link over seed-supplied data are the mitigations; no Sybil defense~\citep{douceur2002sybil} is designed beyond the publishing cost of \S\ref{sec:catalog}.

\subsection{Private search at scale}

On-device search stops scaling when the catalog no longer fits on a phone. The next step in the design is private information retrieval (PIR)~\citep{chor1995pir} against a periodic snapshot of the header catalog, so that a replica answers a query without learning which records were fetched. Single-server lattice-based schemes such as SimplePIR~\citep{henzinger2023simplepir} make this plausible for databases of the relevant size. The scheme is not designed in detail (\S\ref{sec:future}). Until then, a client that fetches individual records from a replica reveals its interest in those records to that replica.
